% Generated draft: requires mathematical and source review.
\documentclass[11pt]{article}
\usepackage[margin=25mm]{geometry}
\usepackage{fontspec}
\setmainfont{Latin Modern Roman}
\usepackage{amsmath,amssymb,mathtools,unicode-math}
\setmathfont{Latin Modern Math}
\usepackage{enumitem,needspace,xurl,hyperref}
\setlength{\parindent}{0pt}
\setlength{\parskip}{5pt}
\setlength{\emergencystretch}{3em}
\newcommand{\N}{\mathbb{N}}
\newcommand{\Z}{\mathbb{Z}}
\newcommand{\Q}{\mathbb{Q}}
\newcommand{\R}{\mathbb{R}}
\newcommand{\C}{\mathbb{C}}
\newcommand{\F}{\mathbb{F}}
\newcommand{\A}{\mathbb{A}}
\newcommand{\PP}{\mathbb{P}}
\newcommand{\E}{\mathbb{E}}
\newcommand{\eps}{\varepsilon}
\newcommand{\dd}{\,\mathrm{d}}
\begin{document}
\section*{UnsolvedMath ID: 30003022}
\textbf{Source problem code:} catalogue entry
\subsection{Definitions and mathematical statement}
Let \(K\) be a field. A valuation is a surjective map \(v:K^{\times}\to\Gamma\) to an ordered abelian group, extended by \(v(0)=\infty\), satisfying \(v(xy)=v(x)+v(y)\) and \(v(x+y)\geq\min\{v(x),v(y)\}\). Write
\[
\mathcal O_v=\{x\in K:v(x)\geq0\},\qquad
\mathfrak m_v=\{x\in K:v(x)>0\},\qquad
Kv=\mathcal O_v/\mathfrak m_v.
\]
The valuation is nontrivial when \(\mathcal O_v\ne K\). It is henselian if, for every monic polynomial \(f\in\mathcal O_v[X]\), each simple root of its reduction \(\overline f\in Kv[X]\) lifts to a root of \(f\) in \(\mathcal O_v\).

A complete first-order theory \(T\) is NIP if no formula \(\varphi(x;y)\) has the independence property: there is no such formula for which, for every positive integer \(n\), some model of \(T\) contains tuples \(a_1,\ldots,a_n\) and \(b_S\), indexed by all subsets \(S\subseteq\{1,\ldots,n\}\), satisfying
\[
\varphi(a_i;b_S)\quad\Longleftrightarrow\quad i\in S.
\]
For the pure-field reading, an NIP field means that \(\operatorname{Th}(K;0,1,+,-,\cdot)\) is NIP. For the valued-structure reading, a valued NIP field means that the theory of the expansion of this structure by a predicate for \(\mathcal O_v\) is NIP. A field is separably closed if it has no proper separable algebraic extension. A field is real closed if it admits an ordering under which every positive element is a square and every polynomial of odd degree has a root.

\begin{enumerate}
\item Under the valued-structure reading, is it true that for every field \(K\) and every valuation \(v\) on \(K\),
\[
\operatorname{Th}(K;0,1,+,-,\cdot,\mathcal O_v)\text{ is NIP}
\quad\Longrightarrow\quad v\text{ is henselian}?
\]
The source does not require \(v\) to be nontrivial in this question. The phrase ``valued NIP field'' is not explicitly disambiguated in the supplied record; if it means only that the underlying field is NIP, the antecedent must instead refer to the pure-field theory, giving a different formulation.
\item For every NIP field \(K\) that is neither separably closed nor real closed, does there exist a nontrivial henselian valuation \(v\) whose valuation ring is definable in \(K\), and for which the resulting valued field \((K,\mathcal O_v)\) is NIP?

More explicitly, under the pure-ring-language, with-parameters reading of definability, the requested existence is that of a valuation \(v\), an integer \(m\geq0\), a ring-language formula \(\theta(x;y_1,\ldots,y_m)\), and a tuple \(a\in K^m\), such that
\[
\mathcal O_v=\{x\in K:K\models\theta(x;a)\},\qquad
\mathcal O_v\ne K,
\]
with \(v\) henselian and \(\operatorname{Th}(K;0,1,+,-,\cdot,\mathcal O_v)\) NIP. Under the parameter-free reading, require \(m=0\). The record does not specify the language of definability or whether parameters are allowed, so these readings remain unresolved.
\end{enumerate}

The second question retains the additional requirement concerning the resulting valued field from the supplied original statement; that requirement is omitted from the supplied clean statement. No infinitude or characteristic restriction is supplied. In particular, finite fields are not excluded by the stated hypotheses of the second question, so the intended domain requires review before any missing infinitude hypothesis could be inserted. These are source questions with unresolved qualifications, not assertions of the displayed implications or existence claims.
\subsection{Short English statement}
The source asks whether every NIP valued field is henselian, and whether every NIP field that is neither separably closed nor real closed admits a definable nontrivial henselian valuation yielding an NIP valued field. The meaning of valued NIP, the language and parameter convention for definability, and the absence of an infinitude assumption require clarification. The original statement includes the resulting-valued-field NIP requirement, although the clean statement omits it.
\subsection{Sources}
\begin{itemize}
\item Local UnsolvedMath registry: permanent ID \(30003022\), title ``Henselian Valuations on NIP Fields,'' exact source problem code \texttt{catalogue entry}, in the supplied complete source record. The record assigns proposed year 2016 and supplies no individual proposer attribution.
\item Supplied report citation: \emph{Model Theory: groups, geometry, and combinatorics} (2016), Oberwolfach Reports, pp.\ 5--50, DOI \texttt{10.4171/owr/2016/1}. \url{https://doi.org/10.4171/owr/2016/1}. This is also the supplied literature source labelled ``Model Theory report.'' The background attributes the collection to workshop open-problem discussions documented in Oberwolfach Reports.
\item Supplied literature source: Jahnke and Simon, \emph{NIP henselian fields}. \url{https://ems.press/content/serial-article-files/32802}.
\item Literature assessment dated exactly \texttt{2026-08-21}: the supplied curation labels the questions open, describes classification results for many NIP henselian fields and substantial transfer results, and says that the general conjectures remain unresolved. This is supplied dated curation, not independently verified here. The background explicitly requests independent verification of time-sensitive status and cited claims.
\item The registry labels the statement \texttt{corrected\_verified} and reports that it was checked against the official report. Those are supplied verification claims, not verification performed for this draft; the supplied original and clean statements differ as described above.
\end{itemize}
\end{document}
